% =====================================================================
%  Sezione 2.4 — Il rumore dei landmark
%  Sostituisce la vecchia sezione "Rumore dei landmark e filtraggio
%  temporale". Prepara la sezione sulle grandezze integrali.
% =====================================================================

\section{Il rumore dei \emph{landmark} e i suoi effetti}
\label{sec:rumore}

L'uso di una rappresentazione basata sulla posa introduce una criticità che
attraversa tutto il resto del lavoro: i \emph{landmark} non sono misure dirette
del corpo, ma stime prodotte automaticamente dal \emph{pose estimator}. Come
tali sono soggetti a rumore, a instabilità e a perdita temporanea del
tracciamento, soprattutto in presenza di occlusioni, di variazioni di
illuminazione, di movimenti rapidi o di punti di vista sfavorevoli.

\subsection{Due manifestazioni diverse}

Il rumore si manifesta in due forme che richiedono trattamenti diversi.

La prima è l'\textbf{assenza della posa}: in una frazione dei fotogrammi lo
stimatore non produce alcuna stima. Sul dataset utilizzato riguarda il 4,5\%
dei fotogrammi, e su un archivio esterno registrato in condizioni diverse
arriva al 19\%. Il trattamento è dichiarativo: il fotogramma viene marcato come
privo di posa e le grandezze valgono zero, ma la marcatura è essenziale perché
uno zero non è un valore neutro. Un angolo del tronco pari a zero descrive un
tronco \emph{orizzontale}, e una sequenza di zeri seguita da un valore reale
produce una derivata di ampiezza arbitraria. Per la stessa ragione, lo stato
temporale dell'estrattore viene azzerato a ogni interruzione, in modo che
nessuna grandezza dinamica attraversi il buco.

La seconda forma è il \textbf{rumore sulla posizione stimata}: il
\emph{landmark} c'è, ma la sua coordinata oscilla anche quando il corpo è
fermo. È la forma più insidiosa, perché non è rilevabile da un controllo di
validità e si propaga silenziosamente a tutte le grandezze derivate.

\subsection{Il filtraggio come prima risposta}

La risposta immediata al secondo problema è il filtraggio temporale. Una media
mobile esponenziale è la scelta naturale in un contesto \emph{edge}, perché è
ricorsiva, non richiede di conservare la storia del segnale e si aggiorna a
ogni fotogramma con una sola operazione:

\[
\tilde{x}_t = \alpha\,x_t + (1-\alpha)\,\tilde{x}_{t-1},
\]

dove $x_t$ è il valore stimato corrente e $\alpha$ regola il compromesso fra
reattività e stabilità. Un valore elevato di $\alpha$ segue rapidamente i
cambiamenti ma lascia passare il rumore; un valore basso stabilizza il segnale
ma attenua proprio le variazioni rapide che caratterizzano una caduta.

Va però dichiarato un fatto sull'implementazione. Nella versione iniziale di
questo lavoro il filtro era \textbf{previsto ma non applicato}: il motore
geometrico accettava i parametri $\alpha$ per posizione e velocità e non li
utilizzava, mentre una diversa componente del sistema ne applicava uno proprio.
Le grandezze utilizzate per l'addestramento e quelle prodotte in esercizio
avevano quindi caratteristiche di segnale differenti. Il difetto è stato
individuato in fase di revisione ed è documentato fra i difetti noti.

\subsection{Perché il filtraggio non è la risposta}

La correzione di quel difetto avrebbe reso coerente il sistema, ma non avrebbe
risolto il problema di fondo, per una ragione che riguarda l'ordine delle
operazioni.

Il filtraggio agisce sulla \emph{posizione}; le grandezze che interessano sono
\emph{derivate} della posizione. Filtrare prima di derivare riduce il rumore ma
riduce anche l'ampiezza del segnale, perché la derivazione amplifica
selettivamente le alte frequenze e il filtro le attenua entrambe senza
distinguerle. Il compromesso fra reattività e stabilità non si risolve: si
sposta.

La misura riportata nella Sezione~\ref{sec:integrali} rende il problema
quantitativo. L'autocorrelazione a un passo delle grandezze derivate si colloca
fra $0{,}18$ e $0{,}44$, cioè in un intervallo in cui la componente incorrelata
è dominante. Nessuna scelta di $\alpha$ porta quei valori nell'intervallo delle
grandezze statiche, che si trovano intorno a $0{,}9$.

La conseguenza pratica è stata osservata sul campo: il classificatore assegnava
punteggio elevato anche a corpi completamente immobili, e la correlazione fra il
punteggio prodotto e il movimento effettivamente presente era $+0{,}139$. Il
modello non stava ignorando la dinamica per una scelta propria: \textbf{ciò che
riceveva come dinamica era in gran parte rumore}.

La risposta adottata non è quindi un filtro migliore, ma una formulazione
diversa delle grandezze dinamiche, descritta nella
Sezione~\ref{sec:integrali}: sostituire la derivazione con l'accumulazione su
una finestra mobile.
